\section{Lampiran}

\subsection{Lampiran A: Dokumentasi Github}

\subsubsection{Link Repository}
% Isi link Github repository di sini
% Contoh: https://github.com/username/DO-THE-MATH


\subsubsection{Dokumentasi Commit dan Kontribusi}
% Isi dokumentasi commit dari setiap anggota
% Screenshot commit history jika perlu


\subsection{Lampiran B: Penggunaan AI dalam Proyek}

\subsubsection{AI Tools yang Digunakan}
% Isi AI tools yang digunakan (Gemini, ChatGPT, dll)
% Contoh:
% - Google Gemini untuk ...
% - ChatGPT untuk ...


\subsubsection{Link Dokumentasi Penggunaan AI}
% Isi link ke chat history AI
% Contoh:
% - Gemini 1: https://...
% - ChatGPT: https://...


\subsubsection{Screenshot Penggunaan AI}
% Tambahkan screenshot chat dengan AI
% \begin{figure}[h]
% \centering
% \includegraphics[width=0.8\textwidth]{gambar/ai_screenshot.png}
% \caption{Screenshot penggunaan AI}
% \end{figure}


\subsection{Lampiran C: Kode Program Lengkap}

\subsubsection{File main.py}
% Isi kode lengkap main.py
% \begin{lstlisting}[caption=main.py]
% # Paste kode di sini
% \end{lstlisting}


\subsubsection{File digit\_recognition.py}
% Isi kode lengkap digit_recognition.py


\subsubsection{File gesture\_tracking.py}
% Isi kode lengkap gesture_tracking.py


% Tambahkan file-file lain sesuai kebutuhan


\subsection{Lampiran D: Screenshot Tambahan}

\subsubsection{Screenshot Proses Development}
% Screenshot selama proses development


\subsubsection{Screenshot Testing}
% Screenshot saat testing aplikasi


\subsubsection{Screenshot Error dan Debugging}
% Screenshot error dan cara debugging (opsional)

\newpage